\section{Qué debe determinar el salario mínimo}
\label{sec:criterios}

La revisión bibliográfica de Orozco organiza la discusión en tres decisiones. Primero, qué nivel del mínimo permite proteger los ingresos de los trabajadores. Segundo, cómo actualizarlo para compensar la inflación e incorporar mejoras reales. Tercero, qué condiciones de empleo, formalidad y capacidad empresarial deben considerarse al tomar esas decisiones.\cite{orozcoRevision} La distinción permite evaluar una propuesta de aumento sin suponer que el nivel inicial era adecuado ni que toda mejora real debe repetirse al año siguiente.

\subsection{Nivel, actualización y salvaguardas}

El \textbf{nivel de referencia} requiere examinar las necesidades de los trabajadores y sus familias y la posición del mínimo en la distribución salarial. Una canasta de consumo informa sobre los gastos que debe cubrir el hogar. La relación entre el mínimo y el salario mediano informa sobre la distancia frente al trabajador situado en el centro de la distribución. Son comparaciones diferentes. El presupuesto de una familia depende de su tamaño, de cuántos integrantes trabajan y de los precios locales. Por ello, un salario mínimo individual no puede compararse directamente con una línea de pobreza por persona sin precisar la composición del hogar. La OIT incorpora estas diferencias al analizar la suficiencia de los salarios.\cite{oitNecesidades}

La \textbf{actualización periódica} parte de un nivel ya fijado. Compensar la inflación evita que ese nivel pierda poder de compra por el aumento de los precios. Incorporar productividad permite que los trabajadores participen del mayor valor agregado que la economía produce por trabajador o por hora. Pero una regla de inflación más productividad no determina si el nivel inicial era adecuado. Si el mínimo era insuficiente, la regla no garantiza corregir esa insuficiencia. Si se había elevado mucho frente a la producción por trabajador, la regla tampoco evalúa cómo las empresas están absorbiendo ese cambio.

Las \textbf{salvaguardas} examinan la posibilidad de sostener el nivel y el ajuste propuestos. La proporción de trabajadores alcanzados por el mínimo, la contratación, las horas, la informalidad y el cumplimiento permiten observar dimensiones que no aparecen en la tasa de productividad nacional. La OIT advierte que la incidencia de un incremento depende también de cuántos trabajadores están afectados y de su efecto sobre la masa salarial.\cite{oitFactores} Estas variables deben interpretarse junto con la actividad económica, porque una caída de la contratación no identifica por sí sola una consecuencia del mínimo.

\subsection{Las experiencias de América responden a objetivos distintos}

El Cuadro \ref{tab:reglas} recoge seis mecanismos documentados en países de América. Los casos ilustran las modalidades discutidas por Orozco. No constituyen una clasificación exhaustiva del continente ni permiten establecer qué mecanismo funciona mejor por el número de países que lo adopta. Una regla publicada y el resultado de una negociación son objetos distintos de comparación.

\begin{table}[H]
\centering\small
\caption{Ejemplos de mecanismos de fijación y actualización del mínimo}
\label{tab:reglas}
\renewcommand{\arraystretch}{1.16}
\begin{tabularx}{\textwidth}{@{}>{\raggedright\arraybackslash}p{2.2cm}>{\raggedright\arraybackslash}X>{\raggedright\arraybackslash}X@{}}
\toprule
Caso & Mecanismo documentado & Distinción relevante\\
\midrule
Canadá, ámbito federal & Ajuste anual con el IPC del año calendario anterior. & Protege frente a precios. Su cobertura es federal y puede prevalecer un mínimo provincial superior.\cite{canada}\\[4pt]
Perú & Metodología de referencia que combina inflación esperada y productividad total de los factores. & La fórmula técnica debe distinguirse de la decisión y la oportunidad de los reajustes.\cite{peru}\\[4pt]
Costa Rica, sector privado & Metodología publicada con costo de vida, corrección por inflación y una porción del crecimiento real por habitante. & El PIB por habitante no equivale a productividad por trabajador. Incluye cláusulas de revisión.\cite{costarica}\\[4pt]
México, zona general & Monto Independiente de Recuperación más un factor de aumento. & Separa un componente de recuperación del ajuste porcentual. El esquema de 2026 difiere en la frontera norte.\cite{mexico}\\[4pt]
Colombia & Concertación con varios criterios económicos y sociales. & No existe una regla que imponga sumar mecánicamente inflación y productividad.\cite{decreto2026}\\[4pt]
Estados Unidos, ámbito federal & Modificación legislativa, sin indexación automática. & El mínimo federal es de USD 7,25 por hora desde 2009. Pueden aplicar pisos estatales superiores.\cite{eeuu}\\
\bottomrule
\end{tabularx}
\fuente{Elaboración del CJC a partir de la revisión de Orozco y las fuentes oficiales citadas. Son casos seleccionados. Para Perú se describe la metodología revisada en Moneda 199 (2024). Para Costa Rica se describe la metodología publicada por el Ministerio de Trabajo, no un cálculo del reajuste de 2027.}
\end{table}

Canadá muestra que la indexación con precios puede constituir una regla explícita de protección del salario. El mínimo federal se ajusta anualmente el primero de abril usando el IPC del año anterior. Si el piso provincial o territorial aplicable es superior, prevalece ese valor.\cite{canada} La experiencia permite identificar una modalidad de actualización, pero no demuestra que indexar solamente con inflación sea apropiado en cualquier país o para cualquier nivel inicial.

Perú y Costa Rica permiten distinguir las medidas que suelen agruparse bajo la palabra productividad. La metodología peruana revisada por el Banco Central de Reserva combina inflación esperada y variación de la productividad total de los factores.\cite{peru} La metodología costarricense publicada utiliza una parte del crecimiento del PIB real por habitante y contempla condiciones que llevan a revisar la aplicación automática de sus componentes.\cite{costarica} La PTF mide la parte del crecimiento que no se atribuye a los insumos productivos medidos. El PIB por habitante también depende de cuántas personas trabajan y de las horas que trabajan. Ninguno es intercambiable con el valor agregado por ocupado o por hora que utiliza este informe.

México incorpora una distinción adicional. La resolución para 2026 separa el Monto Independiente de Recuperación (MIR) del aumento porcentual ordinario. En la zona general se agregaron 17,01 pesos diarios al mínimo anterior y se aplicó un factor de 6,5\% sobre la suma. En la frontera norte el ajuste fue de 5\%, sin MIR.\cite{mexico} El caso muestra cómo una decisión de recuperación puede hacerse explícita y diferenciarse por ámbito. No prueba que Colombia enfrente el mismo rezago ni que deba aplicar el mismo porcentaje.

Estados Unidos ilustra el riesgo de depender de cambios legislativos sin una actualización automática. El mínimo federal sigue en USD 7,25 por hora desde julio de 2009. Con inflación positiva, mantener su valor nominal reduce el poder de compra. Sin embargo, varios estados tienen mínimos superiores, por lo que el piso federal no representa la remuneración mínima aplicable a todos los trabajadores estadounidenses.\cite{eeuu} La comparación relevante es entre mecanismos de actualización, sin equiparar ese congelamiento nominal prolongado con un ajuste colombiano que compense la inflación de un año.

\subsection{Qué implica la revisión para Colombia y para el ajuste de 2027}

La revisión comparada respalda una forma de ordenar la decisión, no un porcentaje único. La compensación de precios, la participación en mejoras productivas y la recuperación de un nivel considerado insuficiente persiguen objetivos diferentes. La concertación debería precisar qué parte del incremento propuesto corresponde a cada objetivo y con qué evidencia se justifica. El crecimiento del PIB, el PIB por habitante, la PTF y la productividad laboral no deben contarse como si fueran ganancias independientes que siempre pudieran sumarse.

Para 2027, este informe propone compensar la inflación de 2026 después de las ganancias reales acumuladas del mínimo. Es una recomendación para el próximo ajuste, no una regla permanente que excluya futuras mejoras reales. Tampoco certifica que el nivel actual cubra las necesidades de todos los hogares. Para resolver esa pregunta hacen falta indicadores de suficiencia y distribución salarial que conserven el tamaño del hogar, el número de perceptores de ingresos y la jornada laboral.

El fundamento del techo provisional de 7\% debe buscarse, por tanto, en el diagnóstico colombiano que se desarrolla a continuación. Los casos internacionales muestran que es posible separar los objetivos del ajuste y definir salvaguardas. No convierten el 7\% en un consenso internacional ni sustituyen el análisis de cuánto han crecido el salario real y la productividad en Colombia.
